% MMW-METHOD-CONTRACT: ee95080ed721168d72ddba0989ae7188b6e8baaceecbefd9c5b20fd61ea1aadf
\section{模型的建立与求解}

\subsection{数据预处理与统一计量口径}

工作负载、GPU 中心、网络时延、功率映射、区域逐时数据与储能信息首先被按字段说明进行一致性核验。工作负载共包含 $50000$ 条输入任务记录，均通过数据完整性校验，缺失单元格数量为 $0$。区域、任务类型和小时索引被统一编码，任务类型与单位等效 GPU 功率按照功率映射表连接，区域间通信采用有向单向时延矩阵。逐时电力数据覆盖任务主时域、末端结清时域和终端结算时点，其中第 $2406$ 小时不得安排任务，仅用于能源与储能状态结算。

任务 $i$ 的持续小时数及有效截止时点被定义为
\begin{equation}
h_i=\frac{d_i}{60},\qquad F_i=\min\{f_i^{\mathrm{valid}},2406\}.
\label{eq:duration_deadline}
\end{equation}
若任务在整数时刻 $s$ 开工，则其与小时区间 $[t,t+1)$ 的实际重叠时长和瞬时占用指示量分别为
\begin{align}
\omega_{it}(s)&=\max\{0,\min(t+1,s+h_i)-\max(t,s)\},\\
\chi_{it}(s)&=\mathbf 1\{\omega_{it}(s)>0\}.
\end{align}
其中，$\omega_{it}(s)$ 被用于 GPU-hour 和电量计算，$\chi_{it}(s)$ 被用于 GPU、IT 功率与设施功率的瞬时容量审计。

系统成本、购电碳排放、有效网络时延和新能源利用率统一按
\begin{align}
C_{\mathrm{op}}
&=\sum_{r,t}\left(c_{rt}^{buy}G_{rt}^{buy}
-c_{rt}^{sell}G_{rt}^{sell}\right)\Delta t,\\
E_{\mathrm{CO2}}
&=\sum_{r,t}\kappa_{rt}G_{rt}^{buy}\Delta t,\\
L_{\mathrm{net}}
&=\frac{\sum_i g_ih_i\ell_{o_ir_i}}{\sum_i g_ih_i},\\
U_{\mathrm{RE}}
&=\frac{\sum_{r,t}(R_{rt}^{use}+R_{rt}^{ch}+R_{rt}^{sell})\Delta t}
{\sum_{r,t}R_{rt}^{av}\Delta t}
\end{align}
复算。网络时延和新能源利用率仅在分母为正时计算，否则相应指标被标记为不可用并退出评分。多目标候选只有通过全部硬约束后才参与比较，其指标差异采用有限正中位绝对尺度标准化\cite{huber1964robust,miettinen1999nonlinear}。

% MMW-ALGORITHM: 固定Huber滞后回归；50000任务NumPy区间切片确定性贪心调度；按真实调度重建AI小时电量和设施负荷；问题二从真实附件中定位区域级MaxGridImport_MW和MaxGridExport_MW并执行购售电边界及互斥审计；问题三使用Baseline AI负荷执行无储能比较与规则型SOC前向策略；问题四使用调度派生FacilityLoad执行联合储能前推并复算情景；MMW_PILOT只读取真实输入并执行有限样本门禁

实际程序采用 Python 与 scikit-learn 实现。需求预测采用固定 Huber 滞后回归，任务调度采用 NumPy 区间切片支持的确定性贪心构造，能源流根据真实调度形成的小时电量重新计算。问题三采用规则型 SOC 前向策略，问题四依据调度派生设施负荷复算各情景；除新能源波动情景采用固定随机种子 $2026$ 外，其余过程均为确定性构造。

\subsection{问题一：需求预测与基础算力调度}

\subsubsection{GPU 需求统计与预测模型}

% MMW-ID: OBJ-Q1
区域 $r$、任务类型 $k$ 在小时 $t$ 的到达 GPU 需求被聚合为
\begin{equation}
D_{rkt}=\sum_{i:o_i=r,\;k_i=k,\;a_i=t}g_i.
\end{equation}
六区域与三类任务的统计结果如表~\ref{tab:q1_group_statistics} 所示。

\begin{table}[htbp]
\centering
\caption{六区域三类任务的需求统计}
\label{tab:q1_group_statistics}
\begin{tabular}{llrrr}
\toprule
来源区域 & 任务类型 & 任务数 & GPU 需求 & GPU-hour\\
\midrule
RegionA & AITraining & 879 & 63070 & 214319.2833\\
RegionA & BatchInference & 3331 & 45066 & 154341.4500\\
RegionA & RealTimeInference & 5852 & 23585 & 79791.7167\\
RegionB & AITraining & 871 & 62551 & 212294.6333\\
RegionB & BatchInference & 3312 & 44274 & 151260.4333\\
RegionB & RealTimeInference & 5017 & 19989 & 67203.0833\\
RegionC & AITraining & 815 & 58899 & 198331.9667\\
RegionC & BatchInference & 2523 & 34710 & 117897.3667\\
RegionC & RealTimeInference & 4221 & 17015 & 58704.9500\\
RegionD & AITraining & 5726 & 407683 & 1387122.6833\\
RegionD & BatchInference & 3346 & 45291 & 155177.2000\\
RegionD & RealTimeInference & 488 & 1993 & 6558.5500\\
RegionE & AITraining & 4091 & 292000 & 978535.1167\\
RegionE & BatchInference & 2172 & 29548 & 100158.3333\\
RegionE & RealTimeInference & 649 & 2628 & 9054.7667\\
RegionF & AITraining & 4177 & 296143 & 1032935.0333\\
RegionF & BatchInference & 2033 & 27173 & 90263.0167\\
RegionF & RealTimeInference & 497 & 1996 & 6837.0667\\
\bottomrule
\end{tabular}
\end{table}

由表~\ref{tab:q1_group_statistics} 可见，各区域的任务组成存在明显差异，因而预测模型需要同时保留区域效应与任务类型效应。RegionD、RegionE 和 RegionF 的 AITraining 任务对应较大的 GPU 需求与 GPU-hour，而 RegionA、RegionB 和 RegionC 的 RealTimeInference 任务数相对较多。任务数、GPU 需求与 GPU-hour 所刻画的侧重点并不相同，因此调度容量不能仅依据任务数量判断，仍须结合持续时间和逐时重叠关系进行审计。

预测模型引入短期、日内和周内滞后项，并采用
\begin{equation}
z_t=\left[
\sin\frac{2\pi t}{24},
\cos\frac{2\pi t}{24},
\sin\frac{2\pi t}{168},
\cos\frac{2\pi t}{168}
\right]^{\mathsf T}
\end{equation}
表示周期特征。未截断预测量为
\begin{equation}
u_{rkt}=\beta_0+
\sum_{\ell\in\{1,2,3,24,48,168\}}\beta_\ell D'_{rk,t-\ell}
+\gamma^{\mathsf T}z_t+\alpha_r+\delta_k,
\end{equation}
最终预测结果取
\begin{equation}
\widehat D_{rkt}=\max\{0,u_{rkt}\}.
\end{equation}

% MMW-ID: CON-Q1-6
预测误差采用
\begin{align}
\mathrm{MAE}&=\frac{1}{n}\sum_j|y_j-\widehat y_j|,\\
\mathrm{RMSE}&=\sqrt{\frac{1}{n}\sum_j(y_j-\widehat y_j)^2}
\end{align}
评价，WAPE 仅在真实值绝对和为正时计算。最终测试的总体 MAE 为 $24.9733$ GPU，总体 RMSE 为 $35.6945$ GPU；测试区间形成 $432$ 条区域—任务类型逐小时预测记录，需求统计明细形成 $43200$ 行结果。

% MMW-LIMITATION: D2
方法契约原计划比较季节性基线并执行完整超参数网格，但实际运行采用固定的 $\epsilon=1.35$ 和 $\alpha=0.001$，未完成全部候选的验证选模。因此，现有误差仅反映已实现固定模型在保留测试区间上的表现。

\subsubsection{任务候选域与容量约束}

% MMW-ID: CON-Q1-1
任务 $i$ 的合法执行区域—开工时刻集合为
\begin{equation}
\Omega_i=\left\{(r,s):s\ge a_i,\ s+h_i\le F_i,\
\ell_{o_ir}\le\ell_i^{\max},\ \omega_{i,2406}(s)=0\right\}.
\end{equation}
若 $\Omega_i$ 为空，则任务及排除原因被记录，构造过程返回结构化状态。

% MMW-ID: CON-Q1-2
每个任务必须且只能选择一个候选，实时推理任务必须在到达时刻开工：
\begin{align}
\sum_{(r,s)\in\Omega_i}x_{irs}&=1,\\
s_i&=\sum_{(r,s)\in\Omega_i}sx_{irs},\qquad c_i=s_i+h_i,\\
s_i&=a_i,\qquad i\in\mathcal I_{\mathrm{RT}}.
\end{align}

% MMW-ID: CON-Q1-3
区域逐时 GPU、IT 功率和设施功率满足
\begin{align}
\sum_{i,s}g_i\chi_{it}(s)x_{irs}&\le G_r^{\max},\\
L_{rt}^{N}+\sum_{i,s}g_ip_{k_i}^{GPU}\chi_{it}(s)x_{irs}
&\le P_r^{IT,\max},\\
\pi_r\left(L_{rt}^{N}+\sum_{i,s}g_ip_{k_i}^{GPU}
\chi_{it}(s)x_{irs}\right)&\le P_r^{F,\max}.
\end{align}

% MMW-ID: CON-Q1-4
能源结算采用实际重叠时长：
\begin{align}
E_{rt}^{AI}&=\sum_{i,s}g_ip_{k_i}^{GPU}\omega_{it}(s)x_{irs},\\
L_{rt}^{AI,avg}&=\frac{E_{rt}^{AI}}{\Delta t},\\
L_{rt}^{F,avg}&=\pi_r(L_{rt}^{N}+L_{rt}^{AI,avg}).
\end{align}

% MMW-ID: CON-Q1-5
区域逐时峰值利用率和 GPU-hour 平均利用率分别为
\begin{align}
U_{rt}^{peak}
&=\frac{\sum_{i,s}g_i\chi_{it}(s)x_{irs}}{G_r^{\max}},\\
U_{rt}^{avg}
&=\frac{\sum_{i,s}g_i\omega_{it}(s)x_{irs}}
{G_r^{\max}\Delta t}.
\end{align}
两项指标仅在 $G_r^{\max}>0$ 时计算。最后 $24$ 小时的区域 GPU 利用情况见图~\ref{fig:gpu_utilization}。

\begin{figure}[htbp]
\centering
\includegraphics[width=0.88\textwidth]{fig_1_gpu_utilization.png}
\caption{最后 $24$ 小时各区域 GPU 利用率}
\label{fig:gpu_utilization}
\end{figure}

图~\ref{fig:gpu_utilization} 展示了统一时间轴上的区域资源占用变化，可用于识别任务集中到达形成的短时波动。峰值利用率反映任何正重叠任务造成的瞬时容量占用，而平均利用率反映按实际重叠时长折算的 GPU-hour，两者不能相互替代。最终容量合规性仍由逐任务区间重叠和瞬时 GPU、IT 与设施功率约束共同判定。

\subsubsection{实际调度算法与代表性结果}

任务按实时属性、有效截止、持续时间、GPU 需求和任务编号形成固定词典序，再在完整候选域中按等待时间、单向时延、加入任务后的峰值 GPU 利用率、区域及开工时刻依次尝试。每次放置后，GPU、IT 功率和设施功率约束均被重新审计。最后 $24$ 小时共有 $538$ 个实际到达任务获得执行区域及开工、完成时刻，其中 $68$ 个任务在第 $2400$ 至第 $2405$ 小时完成结清。表~\ref{tab:q1_schedule_examples} 给出了任务级代表性调度记录，而非完整甘特图。

\begin{table}[htbp]
\centering
\caption{最后 $24$ 小时任务级代表性调度记录}
\label{tab:q1_schedule_examples}
\begin{tabular}{rllllrrrrr}
\toprule
任务ID & 类型 & 来源 & 执行区 & 到达 & 开工 & 完成 & 等待 & 时延 & GPU\\
\midrule
3498 & RealTimeInference & RegionB & RegionB & 2376 & 2376 & 2377.1833 & 0 & 5.0 & 2.0\\
3809 & AITraining & RegionF & RegionF & 2376 & 2376 & 2377.6333 & 0 & 5.0 & 32.0\\
7403 & BatchInference & RegionD & RegionD & 2376 & 2376 & 2382.2333 & 0 & 5.0 & 10.0\\
7902 & BatchInference & RegionE & RegionE & 2376 & 2376 & 2376.4333 & 0 & 5.0 & 11.0\\
9267 & BatchInference & RegionE & RegionE & 2376 & 2376 & 2376.3667 & 0 & 5.0 & 7.0\\
9415 & AITraining & RegionD & RegionD & 2376 & 2376 & 2377.4667 & 0 & 5.0 & 119.0\\
11961 & AITraining & RegionE & RegionE & 2376 & 2376 & 2380.1000 & 0 & 5.0 & 111.0\\
14647 & BatchInference & RegionA & RegionA & 2376 & 2376 & 2377.2333 & 0 & 5.0 & 11.0\\
17820 & RealTimeInference & RegionC & RegionC & 2376 & 2376 & 2377.2500 & 0 & 5.0 & 1.0\\
18201 & BatchInference & RegionE & RegionE & 2376 & 2376 & 2376.6667 & 0 & 5.0 & 16.0\\
26377 & BatchInference & RegionB & RegionB & 2376 & 2376 & 2382.4333 & 0 & 5.0 & 17.0\\
29779 & BatchInference & RegionF & RegionF & 2376 & 2376 & 2381.6000 & 0 & 5.0 & 17.0\\
\bottomrule
\end{tabular}
\end{table}

表~\ref{tab:q1_schedule_examples} 中的记录均在第 $2376$ 小时到达并开工，但任务持续时间和 GPU 需求存在明显差异，说明小时级到达索引不能替代连续完成时刻。表中既包含实时推理，也包含批量推理和 AI 训练任务，可用于核对即时开工、连续占用和任务类型映射。该表仅用于展示任务级字段和时间口径，完整 $50000$ 行任务轨迹见 \texttt{q1\_q2\_q4\_task\_schedule.csv}，不得将该代表性样例解释为完整调度结果。

% MMW-ID: CON-Q1-7
理论模型规定候选耗尽时在局部冲突集中执行不超过 $N_{\mathrm{back}}=2000$ 个节点的有限回溯，并在搜索截止前仍未形成完整方案时返回 \texttt{feasibility\_unresolved}。实际实现未执行该回溯，只保留确定性贪心可行构造。

% MMW-LIMITATION: D1
实际任务调度未实现有限冲突回溯，其探索范围小于方法契约规定的范围。

% MMW-LIMITATION: D8
正式任务构造采用内部时间截止，并为情景复算、审计和文件输出预留运行时间；截止时未完成的部分方案不得作为正式结果。

% MMW-LIMITATION: D10
试跑阶段只读取真实输入并检查有限任务样本的合法时间窗，因此试跑通过不能替代正式全时域可行性审计。

\subsection{问题二：无储能碳感知任务调度}

% MMW-ID: OBJ-Q2
问题二从问题一的可行任务方案出发，在理论模型中以单任务迁移或开工时刻变更构造候选，并对成本、碳排放、有效网络时延、新能源未利用率和等待量进行稳健尺度标准化。实际程序未执行该局部改进，而是复算确定性任务方案。

% MMW-ID: CON-Q2-1
全部 $50000$ 个实际任务均执行唯一选址、实时即时开工、到达与截止、单向时延、瞬时 GPU、IT 功率、设施功率和真实重叠电量约束，预测需求不替代真实任务。

% MMW-ID: CON-Q2-2
问题二不配置储能，故
\begin{equation}
R_{rt}^{ch}=G_{rt}^{ch}=C_{rt}=D_{rt}=0,
\end{equation}
并满足
\begin{align}
R_{rt}^{use}+R_{rt}^{sell}+R_{rt}^{curt}&=R_{rt}^{av},\\
G_{rt}^{buy}+R_{rt}^{use}&=L_{rt}^{F,avg},\\
G_{rt}^{sell}&=R_{rt}^{sell}.
\end{align}

% MMW-ID: CON-Q2-3
购售电边界满足
\begin{align}
0\le G_{rt}^{buy}&\le I_r^{\max}y_{rt}^{grid},\\
0\le G_{rt}^{sell}&\le X_r^{grid}(1-y_{rt}^{grid}),
\qquad y_{rt}^{grid}\in\{0,1\}.
\end{align}
问题二不读取储能售电限制或 SOC 字段。

% MMW-ID: CON-Q2-4
新能源优先分流规则为
\begin{align}
R_{rt}^{use}&=\min\{R_{rt}^{av},L_{rt}^{F,avg}\},\\
G_{rt}^{buy}&=L_{rt}^{F,avg}-R_{rt}^{use}.
\end{align}
新能源富余时，外送量受 $X_r^{grid}$ 限制，其余部分被计为弃电。

% MMW-ID: CON-Q2-5
问题二运行成本为 $-419658147.7076$ 元，碳排放为 $14.3916$ tCO$_2$。全时域累计购电为 $40.4664$ MWh，累计售电为 $1329591.2956$ MWh。由于成本量纲为购电支出减售电收入，负值表示在给定售电价格及外送边界下累计售电收入超过购电支出，并不表示发电成本为负。

% MMW-ID: CON-Q2-6
有效网络时延为 $5.1594$ ms，新能源利用率为 $0.679078$。任务方案中有 $209$ 个任务的执行区域不同于来源区域，末端结清任务为 $68$ 个；GPU、IT 功率、设施功率、网络时延和完成时限审计均通过。

% MMW-LIMITATION: D3
实际程序未执行契约规定的单任务局部改进，$209$ 个迁移任务来自确定性贪心方案，而不是局部搜索收敛结果。

\subsection{问题三：固定负荷下的储能协同模型}

% MMW-ID: OBJ-Q3
问题三固定采用基准 AI IT 负荷和非 AI IT 负荷，并在理论模型中评价成本、碳排放、非负峰值净购电、绝对爬坡及有效新能源未利用率。无储能状态只作为比较对象，不作为含储能构造的启动判据。

% MMW-ID: CON-Q3-1
区域设施负荷及能源流满足
\begin{align}
L_{rt}^{F,avg}&=\pi_r(L_{rt}^{BAI}+L_{rt}^{N}),\\
R_{rt}^{use}+R_{rt}^{ch}+R_{rt}^{sell}+R_{rt}^{curt}&=R_{rt}^{av},\\
G_{rt}^{load}+R_{rt}^{use}+D_{rt}&=L_{rt}^{F,avg},\\
C_{rt}&=R_{rt}^{ch}+G_{rt}^{ch},\\
G_{rt}^{buy}&=G_{rt}^{load}+G_{rt}^{ch},\qquad
G_{rt}^{sell}=R_{rt}^{sell}.
\end{align}

% MMW-ID: CON-Q3-2
储能状态从第 $0$ 小时运行前的初始状态递推：
\begin{align}
S_{r,-1}&=S_r^0,\\
S_{rt}&=S_{r,t-1}+\eta_r^cC_{rt}\Delta t
-\frac{D_{rt}\Delta t}{\eta_r^d},\\
S_r^{min}&\le S_{rt}\le S_r^{max},\\
S_{r,2406}&\ge S_r^0.
\end{align}

% MMW-ID: CON-Q3-3
充放电功率及购售电满足
\begin{align}
0\le C_{rt}&\le C_r^{\max},&
0\le D_{rt}&\le D_r^{\max},\\
0\le G_{rt}^{buy}&\le I_r^{\max}y_{rt}^{grid},&
0\le G_{rt}^{sell}
&\le\min\{X_r^{grid},X_r^{storage}\}(1-y_{rt}^{grid}).
\end{align}
储能充放电及电网购售电均被设置为互斥行为。

% MMW-ID: CON-Q3-4
无储能状态按 $C_{rt}=D_{rt}=0$ 构造。若其违反进口上限，则只记录 \texttt{no\_storage\_baseline\_infeasible}，含储能方案仍从 $S_{r,-1}=S_r^0$ 独立前推。

% MMW-ID: CON-Q3-5
每次储能动作后均检查终端 SOC 的必要可达性：
\begin{equation}
S_{rt}+\sum_{u=t+1}^{2406}\eta_r^cC_r^{\max}\Delta t\ge S_r^0.
\end{equation}
完整轨迹随后接受购电上限、能源守恒、SOC、互斥关系和终端状态审计。

% MMW-ID: CON-Q3-6
净购电、峰值及绝对爬坡量定义为
\begin{align}
N_{rt}&=G_{rt}^{buy}-G_{rt}^{sell},\\
P_r^{peak}&=\max\{0,\max_tN_{rt}\},\\
V_r&=\sum_{t=1}^{2406}|N_{rt}-N_{r,t-1}|.
\end{align}

% MMW-ID: CON-Q3-7
新能源利用率仅在累计可用新能源为正时计算，否则返回结构化不可用状态并退出候选评分。

实际程序采用规则型 SOC 前向策略。无储能与含储能运行成本分别为 $-419820327.7393$ 元和 $-419516103.8102$ 元，成本增加 $304223.9291$ 元。初始 SOC 合计为 $1594.0000$ MWh，终端 SOC 合计为 $3540.0000$ MWh；累计充电量为 $2076.3155$ MWh，累计放电量为 $0$ MWh，且电网充电量为 $0$。无储能与含储能状态的累计购电量均为 $31.5979$ MWh，而累计售电量分别为 $1330267.7686$ MWh 和 $1328783.5422$ MWh。因此，成本增加来自原可售新能源被转化为储能库存后造成的售电收入减少；购电量保持不变，故两种状态的碳排放均为 $8.8506$ tCO$_2$。相应 SOC 轨迹见图~\ref{fig:storage_soc}。

\begin{figure}[htbp]
\centering
\includegraphics[width=0.88\textwidth]{fig_2_storage_soc.png}
\caption{各区域储能 SOC 轨迹}
\label{fig:storage_soc}
\end{figure}

图~\ref{fig:storage_soc} 展示了规则型策略下储能状态的逐时变化，轨迹应结合初始 SOC、上下限和终端约束判断。累计放电量为 $0$ MWh，说明该规则没有形成依靠放电削峰或电价套利的循环。六区域非负峰值净购电功率之和在储能前后均为 $31.5979$ MW，绝对爬坡量均为 $63.1957$ MW，因而该候选未产生削峰或平滑收益。

% MMW-LIMITATION: D4
问题三未执行契约规定的储能动作候选搜索，成本增加及碳排放、峰值和爬坡量不变只描述当前规则型候选。

\subsection{问题四：算—储—电联合调度与情景复算}

% MMW-ID: OBJ-Q4
联合模型评价运行成本、碳排放、有效网络时延、有效等待损失、新能源未利用率及六区域非负峰值净购电功率。各维度仅在有定义、数值有限且候选之间存在正尺度时进入代理评分。

% MMW-ID: CON-Q4-1
联合设施负荷由真实任务方案重新构造：
\begin{equation}
L_{rt}^{F,avg}=\pi_r\left[
L_{rt}^{N}+\frac{1}{\Delta t}
\sum_{i,s}g_ip_{k_i}^{GPU}\omega_{it}(s)x_{irs}
\right].
\end{equation}
任务侧执行问题一的全部硬约束，能源侧执行问题三的能源守恒、SOC、功率和互斥约束。

% MMW-ID: CON-Q4-2
联合初始方案从 $S_{r,-1}=S_r^0$ 独立前推至第 $2406$ 小时，不因问题三无储能比较状态的可行性而停止；若声明预算内未形成完整候选，则返回 \texttt{feasibility\_unresolved}。

% MMW-ID: CON-Q4-3
实时即时开工率、SLA 满足率、按期完成率和等待质量均先执行分母检查。实际程序采用
\begin{equation}
Q_{\mathrm{service}}=\frac{1}{1+\overline W}
\end{equation}
作为平均等待小时的有限代理，结果为 $1.0$。

% MMW-LIMITATION: D7
服务质量未采用按最大允许等待时间归一化的完整指标，故该结果只能解释为平均等待代理。

% MMW-ID: CON-Q4-4
碳约束、峰谷价差、售电价格和新能源水平按方法契约中的离散集合构造，新能源波动情景采用固定种子 $2026$ 并执行非负截断。

% MMW-ID: CON-Q4-5
理论模型规定联合基准与全部情景共享候选评价上限，各情景从同一联合基线独立重构；实际运行共检查 $15$ 个规则型候选。

% MMW-ID: CON-Q4-6
所有候选在启动前和可中断循环内部检查单调时钟，达到搜索截止后丢弃尚未完成的候选，剩余时间只用于文件输出、完整审计和状态整理。

问题四的全部情景结果见表~\ref{tab:q4_scenarios}。

\begin{table}[htbp]
\centering
\caption{问题四的十五个情景比较结果}
\label{tab:q4_scenarios}
\begin{tabular}{llrrrrrrl}
\toprule
情景 & 水平 & 成本/元 & 碳排放 & 时延 & 服务质量 & 新能源利用率 & 峰值/MW & 状态\\
\midrule
基准 & 1.00 & -419353923.7785 & 14.3916 & 5.1594 & 1.0 & 0.679129 & 40.4664 & 完成\\
碳约束 & 0.00 & -419353923.7785 & 14.3916 & 5.1594 & 1.0 & 0.679129 & 40.4664 & 完成\\
碳约束 & 0.10 & -419353923.7785 & 14.3916 & 5.1594 & 1.0 & 0.679129 & 40.4664 & 启发式内未达到\\
碳约束 & 0.20 & -419353923.7785 & 14.3916 & 5.1594 & 1.0 & 0.679129 & 40.4664 & 启发式内未达到\\
碳约束 & 0.30 & -419353923.7785 & 14.3916 & 5.1594 & 1.0 & 0.679129 & 40.4664 & 启发式内未达到\\
峰谷价差 & 0.75 & -419352397.8360 & 14.3916 & 5.1594 & 1.0 & 0.679129 & 40.4664 & 完成\\
峰谷价差 & 1.00 & -419353923.7785 & 14.3916 & 5.1594 & 1.0 & 0.679129 & 40.4664 & 完成\\
峰谷价差 & 1.25 & -419355449.7210 & 14.3916 & 5.1594 & 1.0 & 0.679129 & 40.4664 & 完成\\
售电机制 & 0.00 & 12327.7690 & 14.3916 & 5.1594 & 1.0 & 0.679129 & 40.4664 & 完成\\
售电机制 & 0.50 & -209670798.0048 & 14.3916 & 5.1594 & 1.0 & 0.679129 & 40.4664 & 完成\\
售电机制 & 1.00 & -419353923.7785 & 14.3916 & 5.1594 & 1.0 & 0.679129 & 40.4664 & 完成\\
新能源水平 & 0.80 & -302537839.3757 & 30759.1226 & 5.1594 & 1.0 & 0.805166 & 597.8863 & 完成\\
新能源水平 & 1.00 & -419353923.7785 & 14.3916 & 5.1594 & 1.0 & 0.679129 & 40.4664 & 完成\\
新能源水平 & 1.20 & -457183971.0773 & 0.0000 & 5.1594 & 1.0 & 0.577105 & 0.0000 & 完成\\
新能源波动 & 2026 & -402529632.4199 & 7923.8416 & 5.1594 & 1.0 & 0.675122 & 935.7721 & 完成\\
\bottomrule
\end{tabular}
\end{table}

表~\ref{tab:q4_scenarios} 表明，所有情景均保持 $209$ 个迁移任务，网络时延和服务质量代理也保持不变，这与情景复算固定联合基准任务方案的实现方式一致。售电价格系数为 $0$ 时运行成本为正，而系数为 $0.5$ 和 $1$ 时运行成本为负，说明当前成本指标对售电收入核算较为敏感。新能源水平系数由 $0.8$ 变为 $1.2$ 时，碳排放分别表现为 $30759.1226$ tCO$_2$ 和 $0$ tCO$_2$，峰值净购电分别为 $597.8863$ MW 和 $0$ MW；新能源波动情景则得到 $7923.8416$ tCO$_2$ 的碳排放和 $935.7721$ MW 的峰值净购电。碳削减水平为 $0.1$、$0.2$ 和 $0.3$ 时，目标均在声明启发式内未达到，该状态不构成数学不可行性的证明。

% MMW-LIMITATION: D5
问题四未执行契约规定的共享交替邻域搜索，每个情景只重建并审计一个规则型候选。

% MMW-LIMITATION: D6
碳约束情景未执行额外的碳定向搜索，未达标状态只能表述为在声明启发式内未达到。

% MMW-LIMITATION: D9
只有新能源波动情景采用固定种子 $2026$ 生成随机扰动，其余情景均为确定性复算。

联合基准运行成本为 $-419353923.7785$ 元，碳排放为 $14.3916$ tCO$_2$，网络时延为 $5.1594$ ms，新能源利用率为 $0.679129$，六区域非负峰值净购电功率之和为 $40.4664$ MW。统一约束审计通过，结果被标记为 \texttt{full\_horizon\_feasible=true}。同时，结果被标记为 \texttt{non\_Pareto=true}，且不存在全局或局部最优性证书，因此只能被视为已声明规则及已检查情景中的可行候选。

\subsection{统一约束审计}

任务合法性按
\begin{align}
a_i&\le s_i,\qquad c_i\le F_i,\\
\ell_{o_ir_i}&\le\ell_i^{\max},\\
i\in\mathcal I_{\mathrm{RT}}&\Rightarrow s_i=a_i,\\
\omega_{i,2406}(s_i)&=0
\end{align}
逐项检查。问题二的能源残差以及问题三、四的能源与 SOC 残差分别定义为
\begin{align}
e_{rt}^{q2}
&=G_{rt}^{buy}+R_{rt}^{av}
-L_{rt}^{F,avg}-G_{rt}^{sell}-R_{rt}^{curt},\\
e_{rt}^{energy}
&=G_{rt}^{buy}+R_{rt}^{av}+D_{rt}
-L_{rt}^{F,avg}-C_{rt}-G_{rt}^{sell}-R_{rt}^{curt},\\
e_{rt}^{SOC}
&=S_{rt}-S_{r,t-1}-\eta_r^cC_{rt}\Delta t
+\frac{D_{rt}\Delta t}{\eta_r^d}.
\end{align}
只有任务唯一执行、实时开工、截止、时延、GPU、IT、设施功率、能源守恒、SOC、终端 SOC、购售电互斥、充放电互斥、外送上限及第 $2406$ 小时禁占用全部通过时，才输出全时域可行状态。

% MMW-LIMITATION: L1
实际方法未提供全局或局部最优性证书，所有结论均限于已构造并通过审计的启发式候选。

% MMW-LIMITATION: L2
任务候选和储能动作的完整局部邻域未被穷举，未检查邻域中可能存在更优可行方案。

% MMW-LIMITATION: L3
碳目标未达到只表示在声明启发式内未达到，不构成数学不可行性的证明。

% MMW-LIMITATION: L4
问题四各情景保持联合基准的任务迁移方案不变，无法反映任务迁移方向对情景变化的全部响应。

% MMW-LIMITATION: L5
储能策略未计入循环退化成本，工程经济性评价仍不完整。

% MMW-LIMITATION: L6
复杂实例可能在内部任务构造截止前返回 \texttt{feasibility\_unresolved}，该状态表示运行预算内未完成构造。

% MMW-LIMITATION: L7
问题二只使用区域级最大购电与最大外送功率，不读取储能售电限制或 SOC 字段。

% MMW-LIMITATION: L8
试跑文件仅证明真实输入基础门禁和有限样本时间窗检查通过，不能替代正式全时域审计。
